In this section, we introduce a novel architecture for scalable neural contextual bandit problems, drawing inspiration from Thompson sampling and recent developments in epistemic neural networks (ENN) and EpiNet \cite{osband2021epistemic}. As discussed in Section \ref{sec:ts}, the objective of a deep learning-based Thompson sampling strategy is to accurately estimate the uncertainty of a prediction, pertaining to a context-action pair, while keeping computational costs to a minimum. A significant drawback of many neural methods for Thompson sampling and UCB is their computationally expensive uncertainty estimation process, which impedes their integration into real-world environments. The primary aim of our proposed architecture and design is to address this very issue, while enhancing model performance.

\subsection{Informative Neural Representation}
A key aspect shared across all neural exploration methods is the generation of effective representations for context-action pairs. Within the contextual bandit framework, there are three main representation factors to consider: action representation, context representation, and the interaction between context and action.

We assume a general and unnormalized feature representation from the environment for both contexts and actions. To distill both context and action feature representation vectors, we initially feed each vector into a linear layer followed by a ReLU activation function. This is subsequently followed by a Layer Normalization (LayerNorm) layer \cite{ba2016layer}. Given that raw features can contain disproportionately large values that are challenging to normalize - particularly in an online bandit setting where data continually stream in - it is crucial to employ layer normalization. This technique ensures smoother gradients, promoting stability and generalizability \cite{xu2019understanding}. More formally, 
\begin{equation}
\begin{split}
    h_{\beta_{\text{context}}}(\psi_{S_t}) &= \text{LayerNorm}(\text{ReLU}(\beta_{\text{context}}^T\psi_{S_t})),\\
    h_{\beta_{\text{action}}}(\phi_A) &= \text{LayerNorm}(\text{ReLU}(\beta_{\text{action}}^T\phi_A)),
\end{split}
\end{equation}
where $\theta_{\text{context}} \in \mathbb{R}^{d_{S} x d_{E}}$ and $\theta_{\text{action}} \in \mathbb{R}^{d_{A} x d_{E}}$ are the parameters for context and action summarization. Note that after summarization through these two layers, the outputs are of the same shape.

Having summarized and normalized representations for both action and context, we now model the interaction between these two entities. Collaborative Filtering \cite{resnick1994grouplens, goldberg1992using, herlocker2000explaining} and Matrix Factorization \cite{koren2009matrix} provide intuitive models for this interaction. Inspired by these methods, we use element-wise multiplication to represent interactions,
\begin{equation}
    I(\psi_{S_t}, \phi_A) = h_{\beta_{\text{context}}}(\psi_{S_t}) \odot h_{\beta_{\text{action}}}(\phi_A).
\end{equation}

Utilizing the above three representations, we concatenate the three vectors to derive

\begin{equation} \label{eq:x}
x_{t, A} = \text{concat}\left(h_{\beta_{\text{context}}}(\psi_{S_t}), h_{\beta_{\text{action}}}(\phi_A), I(\psi_{S_t}, \phi_A)\right).
\end{equation}

Notably, this representation diverges from Neural Collaborative Filtering (NCF) \cite{he2017neural} in two primary ways. Firstly, NCF can only manage id-listed features and it aggregates the concatenation of action and context through multi-layer perceptrons. In contrast, our approach can handle feature vectors without making any preliminary assumptions about the input feature shape or values. Secondly, we use the representations obtained through summarization for uncertainty estimation, as discussed in the following subsection. Moreover, we demonstrate in later sections through empirical studies that the direct application of NCF with uncertainty estimation yields inferior performance compared to our method.

\subsection{Epistemic Neural Recommendation (ENR)}
Using the representation derived above, we can now design a neural network for both point estimate and uncertainty estimation. For efficient epistemic uncertainty estimation using neural networks, we employ EpiNet \cite{osband2021epistemic}, an auxiliary architecture for a general neural network. This architecture leverages the final layer representation of a neural network, along with an epistemic index $z$, to generate a sample from the posterior. EpiNet presents a cost-effective method for constructing ENNs, which has been demonstrated to excel in making joint predictions for classification and supervised learning tasks, while achieving strong performance in synthesized neural logistic bandits with minimal additional computational costs. However, a drawback of only utilizing the last layer representation, particularly in a contextual bandit setting, is the diminished representation power. The interactions between context and action are summarized through multi-layers perceptrons for marginal predictions, and considerable information is lost in the process. Motivated by this, we introduce Epistemic Neural Recommendation (ENR) that generates improved posterior samples via more informative representations.

The remaining architecture of ENR is comprised of three parts: a function $f_{\theta_x}: \mathbb{R}^{3d_E} \mapsto \mathbb{R}$ for marginal prediction, and two functions $g_{\sigma}$ and $g_{\sigma^p}$, both mapping from $\mathbb{R}^{3d_E}$ to $\mathbb{R}^{d_z}$, for uncertainty estimation. Upon extraction of $x_{t, A}$, we employ $f_{\theta_x}$ to make the marginal prediction via $f_{\theta_x}(x_{t,A})$.

In addition to $f_{\theta_x}$, $g_{\sigma}$ and $g_{\sigma^p}$ are independently initialized through Glorot Initialization \cite{glorot2010understanding, bartlett2017spectrally}, and $g_{\sigma^p}$ remains constant throughout its lifetime, thereby providing robust regularization. To initialize this segment of the network, ENR samples an epistemic index $z$ from its prior $P_z$, which could either be a discrete one-hot vector or a standard multivariate Gaussian sample. The sampled $z$ is then concatenated with $x_{t, A}$ and processed through both $g_{\sigma}$ and $g_{\sigma^p}$. Consequently, the uncertainty estimation is $(g_{\sigma}(x_{t, A}, z) + g_{\sigma^p}(x_{t, A}, z))^T z$. Hence the final output of ENR is
\begin{equation} \label{eq:forward}
    f_\theta(\psi_S, \phi_A, z) = f_{\theta_x}(x_{t, A}) + (g_{\sigma}(sg[x_{t, A}], z) + g_{\sigma^p}(sg[x_{t, A}], z))^T z,
\end{equation}
where $x_{t, A}$ is defined in Equation \ref{eq:x}, $\theta = (\theta_x, \sigma, \sigma^p, \beta_{\text{context}}, \beta_{\text{action}})$, and $sg[\cdot]$ stops gradient flow of the parameter within the bracket. Note that, if the reward function is a binary variable, then a sigmoid would then be added to the final output of $f_\theta(\psi_S, \phi_A, z)$ as neural logistic regression is a common setup in RS. 

With the neural network above, an ENR Thompson sampling agent executes action 
\begin{equation}
    A_t = \arg\max_{a \in \mathcal{A}}f_\theta(\psi_{S_t}, \phi_a, z), z \sim P_z.
\end{equation}
After receiving reward $R_{t+1}$, then the agent can update the neural network with 
\begin{equation}
    \theta \leftarrow \theta - \alpha\nabla_{\theta}\sum_{z\in \mathcal{Z}}\mathcal{L}(R_{t+1}, f_\theta(\psi_{S_t}, \phi_{A_t}, z))
\end{equation}
through stochastic gradient descent, where $\mathcal{Z}$ is a set of epistemic indices sampled from $P_z$ and $\alpha$ is the learning rate. In practice, we perform batch updates through sampling from a replay buffer and uses other more advanced optimization strategies such as ADAM \cite{kingma2015adam} for parameter updates. The loss function $\mathcal{L}$ here could be cross-entropy loss for neural logistic contextual bandits and mean squared error loss for contextual bandits with non-binary reward. 

For the full illustration of the architecture, please refer to Figure \ref{fig:enr}. For a detailed algorithm illustration, please refer to Algorithm 1. We would like to also note that this algorithm is a full online bandit algorithm without any prior knowledge or pre-training. This algorithm can also be extended to perform offline pretraining before online interactions to reduce the amount of exploration required. We leave this for future work. 

\begin{figure}[t]
     \centering
     \includegraphics[width=0.53\textwidth]{img/ENR.png}
     \caption{Epistemic Neural Recommendation Architecture}
     \label{fig:enr}
     \vspace{-0.1in}
\end{figure}

\begin{algorithm}[t]
\small
	\caption{Epistemic Neural Recommendation (ENR) Thompson Sampling}\label{alg:enr}
	\begin{algorithmic}[1]
	    \State Initialize $\theta, P_z$.
	    \State Initialize a replay buffer $\mathcal{D}$
		\For{$t = 1, 2, \dots$} \Comment{Progressing over time steps}
    		\State Makes observation $O_{t}$, obtains $\psi_{S_t}$ and computes $R_{t}$.
    		\State Store $(\psi_{S_{t-1}}, \phi_{A_{t-1}}, R_t)$ in $\mathcal{D}$.
    		\State Sample $z \sim P_z$ and compute $f_\theta(\psi_{S_t}, \phi_{A}, z) \forall A$ according to Equation \ref{eq:forward}.
    		\State $A_t = \arg\max_{a \in \mathcal{A}}f_{\theta}(\psi_{S_t}, \phi_a, z)$.
    		\For{$k = 1, 2, \dots, K$} \Comment{Optimization}
    	        \State Sample a data batch $B \sim \mathcal{D}$.
    	        \State Sample a set of epistemic indices $Z \sim P_z$
    	        \State $\theta \leftarrow \theta - \alpha\nabla_{\theta}\sum_{i \in B}\sum_{z\in Z}\mathcal{L}(R_{i}, f_\theta(\phi_{S_{i}}, \phi_{A_{i}}, z))$.
    	    \EndFor
		\EndFor
    \end{algorithmic}
    \vspace{-0.03in}
\end{algorithm}


% We first introduce the backbone of architecture we propose, EpiNet. EpiNet is a strategy of building Epistemic Neural Networks with low cost, which has been proven to work well in joint predictions for classification and supervised learning tasks and synthesized neural logistic bandits \cite{osband2021epistemic} with minimal additional computational cost. An EpiNet first initializes a small neural network $g_\eta$ and its prior network $g_{\eta^p}$, both having an output dimension $D_z$. Prior networks are never updated throughout its lifetime and are proven to be effective for regularization \cite{glorot2010understanding, bartlett2017spectrally}. An EpiNet then samples an epistemic index $z \in \mathbb{R}^{D_z}$, sampled its prior $P_z$. Keeping the main neural network architecture unchanged, in the context of contextual bandit, EpiNet first extracts the last layer representation $\sigma_{\beta}(x_{t, A}) \in \mathbb{R}^d$ from the neural network forward pass, where $x_{t, A} = \text{concat}(\phi_{C_t}, \phi_A)$. The EpiNet then concatenates the last layer representation and the epistemic index and sends the concatenated vector through $g_{\eta}$ and $g_{\eta^p}$. The outputs of the two networks are then summed and dot producted with $z$. More formally, the prediction of an EpiNet is
% \begin{equation}
%     h_\theta(x_{t, A}, z) = f_\beta(x_{t, A}) + \left(g_\eta(sg[\sigma_\beta(x_{t, A})], z) + \alpha g_{\eta^p}(sg[\sigma_\beta(x_{t, A})], z)\right)^Tz,
% \end{equation}
% where $f_\beta$ is the main neural network and $sg[\cdot]$ stops the gradient flow from the computation. 
